\chapter{Le pétrole et le gaz naturel}

\noindent\textit{Le pétrole et le gaz naturel sont des ressources fossiles formées sur des
millions d'années. Leur exploitation fournit carburants, matières premières et énergie,
mais leur combustion pose aussi un problème majeur de pollution et de protection de
l'environnement.}
\vspace{8pt}

\connecteBanner

\section{Origine et nature du pétrole}

\begin{definition}[Combustible fossile]
Le \textbf{pétrole} et le \textbf{gaz naturel} sont des \textbf{combustibles fossiles} :
ils se sont formés en plusieurs millions d'années par transformation, sous haute
pression et haute température, de matières organiques (restes d'organismes vivants,
essentiellement marins) accumulées au fond des océans anciens.
\end{definition}

\begin{propriete}[Composition du pétrole brut]
Le pétrole brut est un \textbf{mélange complexe d'hydrocarbures} (molécules formées
uniquement de carbone et d'hydrogène), de tailles de molécules très différentes. Cette
diversité permet, par distillation, d'en extraire des produits aux usages variés.
\end{propriete}

\section{La distillation fractionnée}

\begin{definition}[Distillation fractionnée]
La \textbf{distillation fractionnée} sépare les constituants du pétrole brut selon leur
\textbf{température d'ébullition} : le pétrole est chauffé dans une colonne de
distillation ; les molécules les plus légères (petite taille, température d'ébullition
basse) montent en haut de la colonne, tandis que les molécules les plus lourdes (grande
taille, température d'ébullition élevée) restent en bas.
\end{definition}

\begin{illus}
\begin{tikzpicture}[scale=0.85]
\draw[onbline,line width=1.3pt,fill=onbsoft] (0,0) -- (-0.5,0.3) -- (-0.5,6) -- (0.5,6) -- (0.5,0.3) -- cycle;
\node[right] at (0.7,5.6) {\small gaz (butane, propane)};
\node[right] at (0.7,4.6) {\small essence};
\node[right] at (0.7,3.6) {\small kérosène};
\node[right] at (0.7,2.6) {\small gasoil};
\node[right] at (0.7,1.6) {\small fioul lourd};
\node[right] at (0.7,0.5) {\small bitume};
\draw[->,onbo,line width=1.2pt] (0,-0.4) -- (0,0.2);
\node[below] at (0,-0.5) {\small pétrole brut chauffé};
\node[left,rotate=90] at (-0.9,3) {\small température croissante};
\end{tikzpicture}
\fig{Figure — Colonne de distillation : les fractions se séparent selon leur température d'ébullition.}
\end{illus}

\begin{propriete}[Principales fractions, de la plus légère à la plus lourde]
\textbf{Gaz} (butane, propane) $\to$ \textbf{essence} $\to$ \textbf{kérosène}
(carburant d'avion) $\to$ \textbf{gasoil} $\to$ \textbf{fioul lourd} $\to$
\textbf{bitume} (résidu solide, revêtement routier).
\end{propriete}

\begin{remarque}
Plus une molécule d'hydrocarbure est grande (plus elle contient d'atomes de carbone),
plus sa température d'ébullition est élevée, et plus elle est visqueuse.
\end{remarque}

\section{Combustion des hydrocarbures}

\begin{propriete}[Combustion complète]
En présence de \textbf{quantité suffisante de dioxygène}, un hydrocarbure brûle
complètement et produit du \textbf{dioxyde de carbone} et de l'\textbf{eau} :
\[
\mathrm{CH_4 + 2\,O_2 \longrightarrow CO_2 + 2\,H_2O} \qquad\text{(exemple : méthane).}
\]
\end{propriete}

\begin{exemple}
Combustion complète du butane : $\mathrm{2\,C_4H_{10} + 13\,O_2 \longrightarrow
8\,CO_2 + 10\,H_2O}$.
\end{exemple}

\begin{propriete}[Combustion incomplète]
En \textbf{quantité insuffisante} de dioxygène, la combustion est incomplète et produit
du \textbf{monoxyde de carbone} (CO), gaz \textbf{toxique et incolore}, en plus de l'eau :
\[
\mathrm{2\,CH_4 + 3\,O_2 \longrightarrow 2\,CO + 4\,H_2O}.
\]
\end{propriete}

\begin{remarque}
C'est pourquoi il ne faut jamais utiliser un appareil à combustion (groupe électrogène,
brasero) dans un espace mal aéré : le monoxyde de carbone produit peut provoquer une
intoxication mortelle, sans odeur pour alerter.
\end{remarque}

\section{Pollution et protection de l'environnement}

\begin{propriete}[Impacts environnementaux]
La combustion des hydrocarbures libère du \textbf{dioxyde de carbone} ($\mathrm{CO_2}$),
un \textbf{gaz à effet de serre} qui contribue au réchauffement climatique, ainsi que des
\textbf{particules et polluants} (oxydes d'azote, suies) nuisibles à la qualité de l'air.
L'extraction et le transport du pétrole présentent aussi des risques de \textbf{marées
noires} en cas d'accident.
\end{propriete}

\begin{propriete}[Pistes de protection]
Réduire l'impact environnemental des combustibles fossiles passe par : la \textbf{sobriété
énergétique} (réduire la consommation), le développement des \textbf{énergies
renouvelables} (solaire, hydraulique, éolienne), l'amélioration du \textbf{rendement} des
moteurs et appareils, et le traitement des rejets polluants.
\end{propriete}

% ═══════════════════════════════════════════════════════════════
\section{L'essentiel à retenir}

\begin{synthese}
\textbf{Origine.} Pétrole et gaz naturel : combustibles fossiles issus de matières
organiques transformées sur des millions d'années.\\[4pt]
\textbf{Distillation fractionnée.} Sépare les hydrocarbures du pétrole brut selon leur
température d'ébullition (gaz $\to$ essence $\to$ kérosène $\to$ gasoil $\to$ fioul
$\to$ bitume).\\[4pt]
\textbf{Combustion complète.} Hydrocarbure $+$ dioxygène (suffisant) $\to$
$\mathrm{CO_2}+\mathrm{H_2O}$.\\[4pt]
\textbf{Combustion incomplète.} Dioxygène insuffisant $\to$ production de CO (toxique).\\[4pt]
\textbf{Piège fréquent.} Le monoxyde de carbone (CO) est inodore et incolore : le danger
n'est pas perceptible sans détecteur.
\end{synthese}

% ═══════════════════════════════════════════════════════════════
\section{Méthodes \& savoir-faire}

\methode{Ordonner les fractions du pétrole selon leur volatilité}
Se rappeler que les molécules \textbf{les plus légères} montent en \textbf{haut} de la
colonne (température d'ébullition basse) et les plus lourdes restent en bas.
\begin{exemple}
Le gasoil, plus lourd que l'essence, est extrait plus bas dans la colonne que l'essence.
\end{exemple}

\methode{Écrire et vérifier une équation de combustion complète}
Identifier les produits ($\mathrm{CO_2}$ et $\mathrm{H_2O}$), puis équilibrer en comptant
séparément les atomes de C, H et O de chaque côté.
\begin{exemple}
Propane $\mathrm{C_3H_8}$ : $\mathrm{C_3H_8+5\,O_2\longrightarrow3\,CO_2+4\,H_2O}$
(vérifier : C $3=3$, H $8=8$, O $10=10$).
\end{exemple}

\methode{Distinguer combustion complète et incomplète}
Regarder le produit carboné formé : $\mathrm{CO_2}$ (complète, dioxygène suffisant) ou CO
(incomplète, dioxygène insuffisant, danger).
\begin{exemple}
Une flamme jaune et fuligineuse (avec suie) indique souvent une combustion incomplète.
\end{exemple}

% ═══════════════════════════════════════════════════════════════
\section{Exercices d'application}
\begin{multicols}{2}

\rubrique{Origine et distillation}
\exo{}\diff{1} Qu'est-ce qu'un combustible fossile ?

\exo{}\diff{2} Qu'appelle-t-on distillation fractionnée ?

\exo{}\diff{1} Citer, dans l'ordre, trois fractions obtenues par distillation du pétrole.

\exo{}\diff{2} Pourquoi le bitume est-il extrait tout en bas de la colonne de
distillation ?

\rubrique{Combustion}
\exo{}\diff{2} Écrire l'équation de la combustion complète du méthane $\mathrm{CH_4}$.

\exo{}\diff{2} Quels sont les produits d'une combustion complète d'un hydrocarbure ?

\exo{}\diff{2} Quel gaz toxique peut se former lors d'une combustion incomplète ?

\exo{}\diff{3} Pourquoi le monoxyde de carbone est-il particulièrement dangereux ?

\rubrique{Environnement}
\exo{}\diff{2} Citer un gaz à effet de serre produit par la combustion des hydrocarbures.

\exo{}\diff{2} Citer deux énergies renouvelables pouvant remplacer les combustibles
fossiles.

\exo{}\diff{3} Quel risque environnemental accompagne le transport maritime du pétrole
brut ?

% ═══════════════════════════════════════════════════════════════
\end{multicols}

\section{Exercices d'entraînement}
\begin{multicols}{2}

\rubrique{Origine et distillation}
\exo{}\diff{2} En combien de temps environ se forment le pétrole et le gaz naturel ?

\exo{}\diff{2} Quelle propriété physique des hydrocarbures permet de les séparer par
distillation fractionnée ?

\exo{}\diff{3} Pourquoi les molécules légères montent-elles en haut de la colonne de
distillation ?

\rubrique{Combustion}
\exo{}\diff{2} Équilibrer et vérifier l'équation de combustion complète du propane :
$\mathrm{C_3H_8+O_2\longrightarrow CO_2+H_2O}$.

\exo{}\diff{3} Écrire l'équation de combustion incomplète du méthane produisant du
monoxyde de carbone.

\exo{}\diff{2} Quel signe visuel (couleur de flamme, dépôt) peut indiquer une combustion
incomplète ?

\exo{}\diff{3} Pourquoi ne faut-il jamais faire fonctionner un groupe électrogène dans une
pièce fermée ?

\rubrique{Environnement}
\exo{}\diff{3} Expliquer, en une phrase, le lien entre combustion des hydrocarbures et
réchauffement climatique.

\exo{}\diff{2} Citer deux gestes permettant de réduire sa consommation de combustibles
fossiles au quotidien.

\exo{}\diff{3} En quoi le développement des énergies renouvelables contribue-t-il à la
protection de l'environnement ?

% ═══════════════════════════════════════════════════════════════
\end{multicols}

\section{Sujets type BEPC}

\sujetbac[Distillation du pétrole]
\begin{enumerate}[label=\arabic*)]
\item Qu'est-ce que le pétrole brut ?
\item Sur quel principe repose la distillation fractionnée ?
\item Cite, dans l'ordre, les six principales fractions obtenues, de la plus légère à la
plus lourde.
\item Pourquoi le gasoil est-il extrait plus bas que l'essence dans la colonne ?
\end{enumerate}

\sujetbac[Combustion et sécurité]
\begin{enumerate}[label=\arabic*)]
\item Écris l'équation de combustion complète du méthane $\mathrm{CH_4}$.
\item Quels sont les deux produits de cette combustion complète ?
\item Que se passe-t-il si le dioxygène est en quantité insuffisante ?
\item Pourquoi le monoxyde de carbone est-il particulièrement dangereux pour l'être
humain ?
\end{enumerate}

\sujetbac[Pétrole et environnement]
\begin{enumerate}[label=\arabic*)]
\item Cite un gaz à effet de serre produit par la combustion des hydrocarbures.
\item Quel est l'effet de ce gaz sur le climat ?
\item Cite deux énergies renouvelables alternatives aux combustibles fossiles.
\item Propose une mesure concrète pour réduire la pollution liée aux combustibles
fossiles.
\end{enumerate}

% ═══════════════════════════════════════════════════════════════
\section{Corrigés}
\begin{multicols}{2}

\corrige{de l'exercice 1}
Un combustible formé sur des millions d'années à partir de matières organiques
fossilisées (pétrole, gaz naturel, charbon).

\corrige{de l'exercice 2}
La séparation des constituants du pétrole brut selon leur température d'ébullition.

\corrige{de l'exercice 3}
Par exemple : gaz, essence, kérosène (dans l'ordre de la colonne).

\corrige{de l'exercice 4}
Car c'est la fraction la plus lourde, avec la température d'ébullition la plus élevée :
elle reste en bas de la colonne, non vaporisée.

\corrige{de l'exercice 5}
$\mathrm{CH_4+2\,O_2\longrightarrow CO_2+2\,H_2O}$.

\corrige{de l'exercice 6}
Du dioxyde de carbone ($\mathrm{CO_2}$) et de l'eau ($\mathrm{H_2O}$).

\corrige{de l'exercice 7}
Le monoxyde de carbone (CO).

\corrige{de l'exercice 8}
Car il est incolore et inodore : on ne peut pas détecter sa présence sans appareil, alors
qu'il est toxique, voire mortel.

\corrige{de l'exercice 9}
Le dioxyde de carbone ($\mathrm{CO_2}$).

\corrige{de l'exercice 10}
Par exemple : l'énergie solaire et l'énergie hydraulique (ou éolienne).

\corrige{de l'exercice 11}
Le risque de marée noire en cas d'accident (fuite, naufrage de pétrolier).

\corrige{du sujet 1}
1) Un mélange complexe d'hydrocarbures extrait du sous-sol.\\
2) Sur la différence des températures d'ébullition des différents hydrocarbures qui le
composent.\\
3) Gaz, essence, kérosène, gasoil, fioul lourd, bitume.\\
4) Car le gasoil est constitué de molécules plus grandes, à température d'ébullition plus
élevée, donc il reste plus bas dans la colonne.

\corrige{du sujet 2}
1) $\mathrm{CH_4+2\,O_2\longrightarrow CO_2+2\,H_2O}$.\\
2) Le dioxyde de carbone et l'eau.\\
3) La combustion devient incomplète et produit du monoxyde de carbone (CO), toxique.\\
4) Car il est incolore, inodore, et empêche le sang de transporter correctement le
dioxygène, ce qui peut être mortel sans symptôme d'alerte perceptible.

\corrige{du sujet 3}
1) Le dioxyde de carbone ($\mathrm{CO_2}$).\\
2) Il contribue au réchauffement climatique (effet de serre).\\
3) Par exemple : l'énergie solaire et l'énergie hydraulique.\\
4) Par exemple : développer les transports en commun ou favoriser les énergies
renouvelables pour la production d'électricité.

\end{multicols}
\exercicesBanner
